\documentclass[12pt]{article}

% ---------- Margins ----------
\usepackage[margin=1in]{geometry}

% ---------- Font (Closest to Times in pdfLaTeX) ----------
\usepackage{newtxtext}

% ---------- Spacing ----------
\usepackage{setspace}
\doublespacing{}

% ---------- Paragraph Formatting ----------
\setlength{\parindent}{0.5in}
\setlength{\parskip}{0pt}

\begin{document}

% ---------- MLA Heading (Flush Left) ----------
\noindent
Johnny Alejandro Rojas\\
March 9, 2026\\
COMM 135-C\\
Nataliia Irwin, MS

\vspace{1em}

% ---------- Title ----------
\begin{center}
Self-Evaluation - Informative Speech
\end{center}

\subsection*{AI Use Reflection}

Throughout the process of preparing my speech, I approached AI as a personal writing assistant rather than a replacement for my own thinking. My method began with the knowledge I had built through the course—readings from the textbook, discussions with the professor, and my own engagement with the topic. From that foundation, I would bring my ideas to AI and ask it to help me shape a draft that resembled how I would naturally talk about the subject in conversation. The output was never the final product; it was a mirror. I would read through what was generated and identify the parts that did not sound like me, the phrases I would never choose, or the structure that felt too formal for the voice I wanted to carry into the room. I would then rework those sections, often replacing them with expressions I felt were more creative or more personal. The most effective strategy I developed was using AI to produce a first conversational draft, then treating that draft as raw material to be edited and reused in the final piece. This kept the ideas mine, the structure mine, and ultimately the delivery mine. AI simply helped me get the words moving faster so I could spend more time refining my thinking.

\subsection*{Strengths and Weaknesses}

The feedback I received from both my professor and classmates offered a perspective that surprised me in the best way. Before delivering my speech, I had quietly worried that my topic might come across as too theoretical, too distant from the everyday concerns of the audience. Yet the feedback reflected something different: there was genuine engagement with the subject matter, and several classmates noted that the topic sparked curiosity and held their attention. This was a meaningful strength to discover. At the same time, the feedback also pointed to a challenge that comes naturally with deeper topics: they carry a weight of definitions and background that, if not carefully managed, can stretch the boundaries of the time available. This was an area identified for improvement for me. The more I reflect on the comments I received, the more they align with what I sensed in the room during my own delivery. I felt present in my preparation but, at certain moments, slightly disconnected from the audience in front of me. That disconnect was not rooted in a lack of readiness, but rather I reflect on my own mindset before delivering the topic. The feedback confirmed the content was strong and relevant to the audience; what it also gently revealed was that the bridge between preparation and presence is one I am still learning to build.

\subsection*{Progress}
One of the goals I set in my first self assessment was to take a topic that could easily expand into something more general, and deliver it in a way that was condensed, objective, and grounded in research. Looking back, I can say that goal was met, at least partially. The speech did not wander; it held a direction. It drew on investigation rather than opinion, and it offered the audience a clear informative thread to follow. That, for me, was a step forward. What remains a challenge, and what I did not fully achieve, was the depth of connection I wanted to feel with the people listening. I rehearsed the words. I knew the material. But rehearsing content and rehearsing presence are two different things, and I have come to understand that more clearly now. Preparation, I have learned, is not only about what you know, it is about arriving in the right state of mind, carrying the same energy and spirit that made the topic meaningful to you in the first place. That is the challenge that persists, and it is the one I am most excited and committed to working through.


\subsection*{Goals for Growth}
The goal I am setting for my next speech is to become more efficient with language to say the same things I mean to say, but with fewer words and greater impact. This is not about saying less; it is about saying it better. I want each sentence to carry its own weight. More than that, I want to show up as the person who is genuinely excited to share what they have studied. The image I hold in my mind is simple: speaking to a group of people I care about, friends, family, cherished listeners, and letting a natural development of ideas shape the way I deliver every idea. To work toward this, I plan to rehearse not only the content but the feeling I want to bring into the room. I will practice reading my speech aloud and cutting any phrase that slows the energy without adding meaning. This goal matters because communication is ultimately about connection, and the speaker who connects is not always the most polished one, it is the one who makes the audience feel that what is being shared truly matters.

\end{document}
